\documentclass[12pt,letterpaper]{article}

\usepackage[spanish,es-tabla]{babel}
\usepackage[utf8]{inputenc}
\usepackage[T1]{fontenc}
\usepackage{geometry}
\usepackage{array}
\usepackage{booktabs}
\usepackage{longtable}
\usepackage{listings}
\usepackage{xcolor}
\usepackage{hyperref}
\usepackage{fancyhdr}
\usepackage{graphicx}
\usepackage{float}

\geometry{margin=2.3cm}
\hypersetup{
    colorlinks=true,
    linkcolor=black,
    urlcolor=blue,
    hypertexnames=false,
    pdftitle={Documentacion Proyecto 3 - Simulador AWS S3},
    pdfauthor={Nicole Parra, Alfredo Mercado}
}

\pagestyle{fancy}
\fancyhf{}
\lhead{Proyecto 3 - AWS S3}
\rhead{Principios de Sistemas Operativos}
\cfoot{\thepage}
\setlength{\headheight}{15pt}

\setlength{\parindent}{0pt}
\setlength{\parskip}{6pt}
\renewcommand{\arraystretch}{1.18}

\definecolor{codebg}{RGB}{248,248,248}
\definecolor{codeframe}{RGB}{210,210,210}
\lstdefinestyle{console}{
    basicstyle=\ttfamily\small,
    backgroundcolor=\color{codebg},
    frame=single,
    rulecolor=\color{codeframe},
    breaklines=true,
    columns=fullflexible,
    keepspaces=true,
    showstringspaces=false
}
\lstdefinestyle{cstyle}{
    language=C,
    basicstyle=\ttfamily\small,
    backgroundcolor=\color{codebg},
    frame=single,
    rulecolor=\color{codeframe},
    breaklines=true,
    columns=fullflexible,
    keepspaces=true,
    showstringspaces=false
}

\begin{document}

\begin{titlepage}
    \centering
    \vspace*{1cm}
    {\Large Tecnologico de Costa Rica\par}
    {\Large Escuela de Computacion\par}
    {\Large Principios de Sistemas Operativos\par}
    \vspace{2.2cm}
    {\Huge\bfseries Tercer Proyecto\par}
    \vspace{0.5cm}
    {\LARGE Simulador de AWS S3 sobre sockets\par}
    \vspace{2cm}
    {\large Profesor: Armando Arce Orozco\par}
    \vspace{0.8cm}
    {\large Estudiantes:\par}
    {\large Nicole Parra, 2023223291\par}
    {\large Alfredo Mercado, 2023377984\par}
    \vfill
    {\large Junio 2026\par}
\end{titlepage}

\pagenumbering{roman}
\tableofcontents
\newpage
\pagenumbering{arabic}

\section{Introduccion}

El proyecto implementa una simulacion de almacenamiento remoto de objetos al estilo
AWS S3. La solucion contiene dos ejecutables principales: \texttt{aws-s3\_server},
que administra los buckets en el servidor, y \texttt{aws-s3}, que funciona como
cliente de linea de comandos. Ambos programas estan escritos en C y se comunican
por sockets TCP.

La idea central es representar cada bucket como un unico archivo local dentro del
servidor. En ese archivo se guarda primero un bloque de directorio con metadatos y,
despues de ese bloque, el contenido binario de los objetos. Las carpetas no existen
como estructuras fisicas; se simulan mediante prefijos en las claves, igual que en S3.
Por ejemplo, \texttt{fotos/2026/playa.jpg} es una sola clave con barras diagonales,
no una jerarquia real de directorios dentro del bucket.

\section{Descripcion del problema}

El objetivo es construir un sistema cliente-servidor que permita gestionar archivos
locales y remotos usando comandos similares a \texttt{aws s3}. El cliente debe
recibir los argumentos desde terminal, interpretarlos como rutas locales o rutas
\texttt{s3://bucket/clave}, y enviar la operacion correspondiente al servidor.

El servidor debe mantener multiples buckets. Cada bucket se guarda en un unico
archivo, con un directorio interno que registra las claves de los objetos, el offset
donde inicia cada contenido y el tama\~{n}o en bytes. Este requisito evita guardar cada
objeto como un archivo separado y obliga a administrar espacio libre dentro del
archivo del bucket.

Cuando un objeto se reemplaza, hay dos casos:
\begin{itemize}
    \item Si el nuevo archivo tiene el mismo tama\~{n}o, se sobrescribe en el mismo offset.
    \item Si el tama\~{n}o cambia, el espacio anterior se marca como libre y el nuevo
    contenido se escribe en un hueco disponible o al final del bucket.
\end{itemize}

Ademas, al eliminar objetos se debe actualizar la tabla de objetos y agregar el
espacio que ocupaban a una lista de huecos reutilizables. La solucion implementada
usa primer ajuste para asignar espacio libre: se toma el primer hueco suficientemente
grande y, si sobra espacio, el hueco se reduce.

\section{Comandos implementados}

\begin{longtable}{>{\ttfamily}p{2.2cm}p{10.7cm}}
\toprule
\textnormal{\textbf{Comando}} & \textbf{Descripcion} \\
\midrule
\endhead
ls & Lista todos los buckets cuando no recibe argumentos, o lista objetos y prefijos dentro de un bucket. \\
mb & Crea un bucket nuevo usando una ruta \texttt{s3://bucket}. \\
cp & Copia archivos local a S3, S3 a local o S3 a S3. Soporta \texttt{--recursive} para directorios o prefijos. \\
mv & Mueve un archivo u objeto mediante copia seguida de eliminacion del origen. Tambien soporta \texttt{--recursive}. \\
rm & Elimina un objeto. Con \texttt{--recursive} elimina todos los objetos que empiezan con el prefijo indicado. \\
sync & Sincroniza local a S3 o S3 a local. Con \texttt{--delete} elimina en el destino lo que ya no existe en el origen. \\
rb & Elimina un bucket vacio. Con \texttt{--force} permite eliminarlo aunque tenga objetos. \\
\bottomrule
\end{longtable}

\section{Definicion de estructuras de datos utilizadas}

Las estructuras principales estan definidas en \texttt{include/common.h} y
\texttt{include/bucket.h}. Se eligieron estructuras simples, con tama\~{n}os maximos
fijos, porque el enfoque es la administracion de buckets y no una base de datos dinamica de metadatos.

\subsection{URI S3}

\begin{lstlisting}[style=cstyle]
typedef struct {
    char bucket[S3_MAX_BUCKET];
    char key[S3_MAX_KEY];
    int is_s3;
} s3_uri_t;
\end{lstlisting}

\texttt{s3\_uri\_t} guarda el resultado de interpretar una ruta
\texttt{s3://bucket/clave}. El campo \texttt{bucket} contiene el nombre del bucket,
\texttt{key} contiene la clave o prefijo dentro del bucket, e \texttt{is\_s3}
indica si el argumento era una ruta remota.

\subsection{Entrada de objeto}

\begin{lstlisting}[style=cstyle]
typedef struct {
    int used;
    char key[S3_MAX_KEY];
    uint64_t offset;
    uint64_t size;
} bucket_entry_t;
\end{lstlisting}

Cada \texttt{bucket\_entry\_t} representa un objeto almacenado. La clave es plana y
puede contener barras para simular carpetas. \texttt{offset} indica el byte dentro
del archivo \texttt{.bucket} donde empieza el contenido, y \texttt{size} indica
cuantos bytes deben leerse o transmitirse.

\subsection{Hueco libre}

\begin{lstlisting}[style=cstyle]
typedef struct {
    int used;
    uint64_t offset;
    uint64_t size;
} bucket_hole_t;
\end{lstlisting}

La lista de \texttt{bucket\_hole\_t} permite reutilizar espacio interno cuando se
borra un objeto o cuando se reemplaza por otro con diferente tama\~{n}o. El algoritmo
implementado es primer ajuste.

\subsection{Directorio del bucket}

\begin{lstlisting}[style=cstyle]
typedef struct {
    char magic[8];
    uint32_t entry_count;
    uint32_t hole_count;
    bucket_entry_t entries[BUCKET_MAX_OBJECTS];
    bucket_hole_t holes[BUCKET_MAX_HOLES];
} bucket_header_t;
\end{lstlisting}

El \texttt{bucket\_header\_t} esta al inicio de cada archivo de bucket. Incluye una
firma \texttt{S3BKT01}, el conteo de objetos activos, el conteo de huecos libres, la
tabla de objetos y la tabla de huecos. En esta version se soportan hasta 512 objetos
y 512 huecos por bucket.

\section{Descripcion detallada de componentes}

\subsection{\texttt{src/common.c}}

Este modulo contiene funciones compartidas por cliente y servidor:
\begin{itemize}
    \item Apertura de sockets cliente y servidor.
    \item Lectura y escritura exacta de bytes con \texttt{read\_exact} y
    \texttt{write\_exact}.
    \item Lectura de lineas de control con \texttt{read\_line}.
    \item Envio de lineas formateadas con \texttt{send\_linef}.
    \item Copia de streams binarios entre archivos y sockets.
    \item Utilidades de rutas, creacion de directorios y calculo de tama\~{n}os.
\end{itemize}

La separacion es importante porque evita duplicar codigo sensible en cliente y
servidor. En especial, las funciones de lectura y escritura exacta protegen contra
lecturas parciales del socket.

\subsection{\texttt{src/bucket.c}}

Este modulo implementa la capa de almacenamiento. Sus responsabilidades son:
\begin{itemize}
    \item Crear y eliminar archivos \texttt{.bucket}.
    \item Leer y escribir el bloque de directorio.
    \item Insertar objetos desde archivos locales o desde streams recibidos por socket.
    \item Enviar objetos al cliente.
    \item Copiar objetos entre buckets.
    \item Eliminar objetos individuales o prefijos completos.
    \item Listar buckets y claves.
    \item Consultar el tama\~{n}o de una clave para \texttt{sync}.
\end{itemize}

El bloque de directorio se reescribe despues de cada operacion que modifica metadatos.
El contenido de los objetos se conserva en la zona de datos. Si se libera espacio, se
agrega un hueco a la tabla para reutilizacion futura.

\subsection{\texttt{src/server.c}}

El servidor crea un socket TCP, escucha en el puerto configurado y acepta conexiones
de clientes. Por cada conexion hace \texttt{fork}, de modo que cada comando se atiende
en un proceso hijo. La conexion procesa una linea de comando, ejecuta la operacion y
responde con \texttt{OK} o \texttt{ERR}.

El protocolo interno usa nombres de comandos mas directos que la interfaz del cliente.
Los comandos internos son \texttt{CREATE\_BUCKET}, \texttt{REMOVE\_BUCKET},
\texttt{LIST}, \texttt{LIST\_RAW}, \texttt{PUT}, \texttt{GET}, \texttt{COPY},
\texttt{DELETE} y \texttt{STAT}.

\subsection{\texttt{src/client.c}}

El cliente implementa la interfaz. Primero procesa opciones
globales como \texttt{--host} y \texttt{--port}; luego identifica el comando y separa
rutas locales de rutas remotas con \texttt{parse\_s3\_uri}.

Para operaciones recursivas, el cliente recorre directorios locales con \texttt{opendir}
y \texttt{readdir}. Para operaciones sobre prefijos remotos, pide al servidor una lista
cruda con \texttt{LIST\_RAW} y luego copia, descarga o elimina cada clave coincidente.

\section{Mecanismo de creacion de archivos}

Los buckets se guardan en el directorio \texttt{buckets/}. Cada bucket se almacena
como un archivo con extension \texttt{.bucket}. La creacion sigue estos pasos:

\begin{enumerate}
    \item El servidor valida el nombre del bucket.
    \item Crea el directorio \texttt{buckets/} si no existe.
    \item Crea el archivo \texttt{.bucket} en modo binario.
    \item Escribe un \texttt{bucket\_header\_t} inicial con la firma
    \texttt{S3BKT01}, sin objetos y sin huecos.
\end{enumerate}

La zona de datos empieza despues del \texttt{bucket\_header\_t}. Cuando se guarda un
objeto nuevo, el servidor busca primero un hueco libre suficiente. Si lo encuentra,
escribe ahi; si no, escribe al final del archivo.

\section{Mecanismo de comunicacion con el servidor}

El protocolo combina lineas de control en texto con payloads binarios. Esto mantiene
los comandos legibles y evita corromper archivos que no son texto.

\subsection{Carga de un archivo}

Para subir un objeto, el cliente envia:

\begin{lstlisting}[style=console]
PUT bucket key size\n
<size bytes exactos del archivo>
\end{lstlisting}

El servidor lee la linea, reserva espacio en el bucket y luego lee exactamente
\texttt{size} bytes desde el socket. Al terminar, responde \texttt{OK} o \texttt{ERR}.

\subsection{Descarga de un archivo}

Para descargar un objeto, el cliente envia:

\begin{lstlisting}[style=console]
GET bucket key\n
\end{lstlisting}

Si el objeto existe, el servidor responde:

\begin{lstlisting}[style=console]
OK size\n
<size bytes exactos del objeto>
\end{lstlisting}

El cliente crea los directorios locales necesarios y escribe los bytes recibidos en el
archivo destino.

\subsection{Listados y sincronizacion}

\texttt{LIST} devuelve una vista amigable con objetos y prefijos. \texttt{LIST\_RAW}
devuelve claves completas y tama\~{n}os, separadas por tabulador. Esta segunda variante se
usa para \texttt{cp --recursive}, \texttt{mv --recursive}, \texttt{rm --recursive} y
\texttt{sync}.

\section{Pruebas de funcionamiento}

El proyecto incluye una prueba automatizada en \texttt{tests/smoke.sh}. Esta prueba
compila el proyecto, levanta un servidor local, ejecuta los comandos principales y usa
\texttt{cmp} para validar que los archivos descargados son byte a byte iguales a los
originales.

\begin{lstlisting}[style=console]
make test
\end{lstlisting}

La secuencia cubierta por la prueba es:

\begin{lstlisting}[style=console]
./aws-s3 --port 19090 mb s3://demo
./aws-s3 --port 19090 cp tmp-smoke/local/a.txt s3://demo/
./aws-s3 --port 19090 cp tmp-smoke/local s3://demo/copia/ --recursive
./aws-s3 --port 19090 ls s3://demo/
./aws-s3 --port 19090 cp s3://demo/a.txt tmp-smoke/out/a.txt
cmp tmp-smoke/local/a.txt tmp-smoke/out/a.txt
./aws-s3 --port 19090 cp s3://demo/copia/ tmp-smoke/out/copia --recursive
cmp tmp-smoke/local/dir/b.txt tmp-smoke/out/copia/dir/b.txt
./aws-s3 --port 19090 cp s3://demo/copia/ s3://demo/copia2/ --recursive
./aws-s3 --port 19090 mv s3://demo/a.txt s3://demo/movido.txt
./aws-s3 --port 19090 mv s3://demo/copia2/ s3://demo/copia3/ --recursive
./aws-s3 --port 19090 cp s3://demo/copia3/dir/b.txt tmp-smoke/out/b2.txt
cmp tmp-smoke/local/dir/b.txt tmp-smoke/out/b2.txt
./aws-s3 --port 19090 rm s3://demo/copia/ --recursive
./aws-s3 --port 19090 sync tmp-smoke/local s3://demo/sync/ --delete
./aws-s3 --port 19090 sync s3://demo/sync/ tmp-smoke/out/sync
cmp tmp-smoke/local/dir/b.txt tmp-smoke/out/sync/dir/b.txt
./aws-s3 --port 19090 rb s3://demo --force
\end{lstlisting}

La salida esperada al final de la prueba es:

\begin{lstlisting}[style=console]
Pruebas smoke completadas
\end{lstlisting}

\section{Pruebas de rendimiento}

Para revisar el comportamiento con archivos de diferentes tama\~{n}os, se ejecuto una
prueba manual de rendimiento en ambiente Unix. La prueba mide carga, descarga e
integridad usando archivos binarios de 10 MiB, 50 MiB y 100 MiB, con tres
repeticiones por tama\~{n}o. La medicion se hizo contra el servidor local por
\textit{loopback}; por eso los resultados representan el costo de protocolo,
copias de bytes y acceso a disco local, no el rendimiento de una red real.

\begin{lstlisting}[style=console]
make
rm -rf buckets tmp-perf
mkdir -p tmp-perf/in tmp-perf/out
dd if=/dev/urandom of=tmp-perf/in/10m.bin bs=1M count=10 status=none
dd if=/dev/urandom of=tmp-perf/in/50m.bin bs=1M count=50 status=none
dd if=/dev/urandom of=tmp-perf/in/100m.bin bs=1M count=100 status=none
./aws-s3_server 19191 > tmp-perf/server.log 2>&1 &
SERVER_PID=$!
sleep 1
./aws-s3 --port 19191 mb s3://perf

echo "archivo,bytes,rep,subida_s,descarga_s,integridad"
# Para cada archivo y repeticion se midio con date +%s%N.
# La integridad se valido con cmp -s.
./aws-s3 --port 19191 rb s3://perf --force
kill "$SERVER_PID"
\end{lstlisting}

Los resultados crudos observados fueron:

\begin{table}[H]
\centering
\small
\begin{tabular}{lrrrrc}
\toprule
Archivo & Bytes & Rep. & Subida (s) & Descarga (s) & Integridad \\
\midrule
10m.bin & 10485760 & 1 & 0.0120 & 0.0109 & OK \\
10m.bin & 10485760 & 2 & 0.0087 & 0.0108 & OK \\
10m.bin & 10485760 & 3 & 0.0091 & 0.0115 & OK \\
50m.bin & 52428800 & 1 & 0.0490 & 0.0462 & OK \\
50m.bin & 52428800 & 2 & 0.0467 & 0.0448 & OK \\
50m.bin & 52428800 & 3 & 0.0341 & 0.0439 & OK \\
100m.bin & 104857600 & 1 & 0.0955 & 0.0927 & OK \\
100m.bin & 104857600 & 2 & 0.0585 & 0.0913 & OK \\
100m.bin & 104857600 & 3 & 0.0670 & 0.0898 & OK \\
\bottomrule
\end{tabular}
\caption{Resultados crudos de carga, descarga y comparacion con \texttt{cmp}.}
\end{table}

Con esos datos se calcularon los promedios y el rendimiento efectivo:

\begin{table}[H]
\centering
\begin{tabular}{lrrrr}
\toprule
Archivo & Tama\~{n}o (MiB) & Subida prom. (s) & Descarga prom. (s) & Integridad \\
\midrule
10m.bin & 10 & 0.0099 & 0.0111 & 3/3 OK \\
50m.bin & 50 & 0.0433 & 0.0450 & 3/3 OK \\
100m.bin & 100 & 0.0737 & 0.0913 & 3/3 OK \\
\bottomrule
\end{tabular}
\caption{Promedio de tres repeticiones por tama\~{n}o.}
\end{table}

\begin{table}[H]
\centering
\begin{tabular}{lrr}
\toprule
Archivo & Subida (MiB/s) & Descarga (MiB/s) \\
\midrule
10m.bin & 1006.7 & 903.6 \\
50m.bin & 1155.6 & 1111.9 \\
100m.bin & 1357.5 & 1095.7 \\
\bottomrule
\end{tabular}
\caption{Rendimiento efectivo sobre loopback local.}
\end{table}

Estas pruebas son importantes porque validan dos aspectos:
\begin{itemize}
    \item El protocolo transmite bytes exactos y no depende de que el archivo sea texto.
    \item El bloque de directorio sigue siendo correcto despues de varias inserciones,
    descargas y eliminaciones.
\end{itemize}

\section{Analisis de resultados}

La prueba automatizada, valida que
un bucket pueda crearse, recibir archivos, listar contenido, copiar prefijos completos,
mover objetos, eliminar prefijos, sincronizar directorios y eliminarse con
\texttt{--force}.

Las comparaciones con \texttt{cmp} son la verificacion mas importante porque comprueban
que la descarga produce exactamente los mismos bytes que el archivo original. Esto
demuestra que la separacion entre linea de control y payload binario funciona
correctamente.

La prueba de rendimiento con archivos de 10 MiB, 50 MiB y 100 MiB mostro tiempos
crecientes conforme aumenta el tama\~{n}o del archivo. En todas las repeticiones,
\texttt{cmp} confirmo que el archivo recuperado era identico al original. El mayor
rendimiento observado en archivos grandes es esperable porque el costo fijo de abrir
conexiones, procesar comandos y actualizar el bloque de directorio se reparte entre
mas bytes transferidos.

La sincronizacion compara existencia y tama\~{n}o de archivos. Esta decision mantiene
el dise\~{n}o simple y suficiente para el alcance del proyecto.

\section{Limitaciones}

\begin{itemize}
    \item Cada bucket soporta hasta 512 objetos y 512 huecos libres.
    \item Las claves tienen un maximo de 255 caracteres.
    \item El protocolo separa argumentos por espacios, por lo que no esta pensado para
    claves remotas con espacios.
    \item \texttt{sync} detecta cambios por tama\~{n}o; no compara hashes ni tiempos de
    modificacion.
    \item El servidor atiende procesos hijos con \texttt{fork}, pero no implementa
    bloqueo de escritura entre operaciones concurrentes sobre el mismo bucket.
\end{itemize}

Estas limitaciones no impiden cumplir los comandos solicitados, pero delimitan el alcance
real de la implementacion.

\section{Compilacion y ejecucion}

El proyecto esta preparado para ambiente Unix. La
compilacion se realiza con:

\begin{lstlisting}[style=console]
make
\end{lstlisting}

Luego se levanta el servidor:

\begin{lstlisting}[style=console]
./aws-s3_server 9090
\end{lstlisting}

Y en otra terminal se ejecutan comandos del cliente:

\begin{lstlisting}[style=console]
./aws-s3 --port 9090 mb s3://demo
./aws-s3 --port 9090 cp archivo.txt s3://demo/
./aws-s3 --port 9090 ls s3://demo/
./aws-s3 --port 9090 cp s3://demo/archivo.txt ./
./aws-s3 --port 9090 rm s3://demo/archivo.txt
./aws-s3 --port 9090 rb s3://demo --force
\end{lstlisting}

Para ejecutar la validacion automatizada:

\begin{lstlisting}[style=console]
make test
\end{lstlisting}

\section{Conclusiones}

La implementacion cumple el objetivo de simular un servicio de almacenamiento de
objetos con buckets, claves planas y prefijos. El cliente traduce comandos similares a
AWS S3 en mensajes simples sobre sockets, mientras que el servidor administra los
metadatos y el contenido dentro de archivos \texttt{.bucket}.

El uso de un bloque de directorio fijo simplifica la recuperacion de metadatos y permite
mantener cada bucket en un solo archivo. La lista de huecos libres resuelve el problema
de reutilizar espacio cuando se eliminan o reemplazan objetos, y el algoritmo de primer
ajuste es suficiente para esta escala.

Las pruebas propuestas cubren los comandos principales y validan la integridad de los
archivos recuperados. Para evitar problemas de ejecucion en la entrega, el proyecto debe
compilarse y probarse en un ambiente Unix con \texttt{gcc}, \texttt{make} y utilidades
POSIX disponibles.

\end{document}
